% ═══════════════════════════════════════════════════════════════════════════════
%   CausalWorlds 2026 Extended Abstract (Grenoble, June 22–26)
%   1–3 pages excluding references
% ═══════════════════════════════════════════════════════════════════════════════
\documentclass[11pt,a4paper]{article}

\usepackage{fontspec}
\defaultfontfeatures{Ligatures=TeX}
\usepackage{amsmath,amssymb}
\usepackage{booktabs}
\usepackage{tikz}
\usetikzlibrary{arrows.meta,positioning}
\usepackage{xcolor}
\usepackage[colorlinks=true,linkcolor=blue!70!black,citecolor=blue!70!black,urlcolor=blue!70!black]{hyperref}

% ─── Page geometry ──────────────────────────────────────────────────────────
\usepackage[margin=2.3cm]{geometry}
\setlength{\parskip}{3pt plus 1pt minus 1pt}
\setlength{\parindent}{0pt}

% ─── Compact sections ──────────────────────────────────────────────────────
\usepackage{titlesec}
\titleformat{\section}{\normalfont\bfseries\large}{\thesection.}{0.5em}{}
\titlespacing*{\section}{0pt}{8pt plus 2pt minus 2pt}{4pt plus 1pt minus 1pt}

% ─── Compact captions ──────────────────────────────────────────────────────
\usepackage[font=small,labelfont=bf,skip=4pt]{caption}

\begin{document}

% ─── Title block ────────────────────────────────────────────────────────────
\begin{center}
  {\Large\bfseries An $O(N)$ Kuratowski Calculus Engine for Causal Sets:\\[2pt]
  Emergent Particle Generations from\\[2pt]
  Topological Defects at $N = 10^6$}

  \medskip
  {\normalsize Juan Pablo Silva Alvarado}\\
  {\small Metric Engineering Division, Bogot\'{a}, Colombia}\\
  {\small\url{https://github.com/ertwro/fractal_entropic_geometrodynamics}}

  \medskip
  {\small\itshape Extended abstract submitted to CausalWorlds 2026 (Grenoble, June 22--26)}
\end{center}

\medskip

% ═══════════════════════════════════════════════════════════════════════════════
%   Abstract
% ═══════════════════════════════════════════════════════════════════════════════

\textbf{Abstract.}
We present an $O(N)$ simulation engine for Poisson-sprinkled causal sets
that detects topological defects---\emph{Causal Prisms} ($K_{2,n}$ bipartite
subgraphs)---in the Hasse diagram and classifies them into exactly three
generations via a topological invariant.
At $N = 10^6$: spectral dimension $d_S = 3.64$ at the causal diamond core
(approaching the 4D continuum); mass hierarchy
Gen\,1\,:\,Gen\,2\,:\,Gen\,3 $= 4.49 : 6.51 : 7.80$;
CPT symmetry between matter and antimatter prisms.
The $O(N)$ scaling is not an engineering optimisation---it is
a computational proof that the physics is local
(bounded Hasse degree $D \leq 15$).
The generation bound $g \leq 3$ follows from the sign trichotomy of bulk
momentum, not from model fitting.
Open-source Rust implementation and pedagogic booklets available at
the repository above.

% ═══════════════════════════════════════════════════════════════════════════════
%   Section 1
% ═══════════════════════════════════════════════════════════════════════════════
\section{The Engine --- $O(N)$ Kuratowski Calculus}

We built a simulation that produces three particle generations, a mass
hierarchy, and CPT symmetry from nothing but Poisson-sprinkled points and
integer combinatorics.  It runs at $O(N)$ and processes $10^6$ events in
63\,seconds on a standard laptop.

The engine pipeline has four phases:

\textbf{Phase 1: Sprinkling \& Hasse diagram}
(\texttt{diamond.rs}).
$N$ points are Poisson-sprinkled into a 4D causal diamond
($V = \pi T^4/24$, $\rho \approx 1$).
The Hasse diagram (transitive reduction of the causal order) is built
via an Occupied Cell Index: for each new event, only cells within the
causal past light cone are probed, and three termination valves
enforce strict locality~\cite{bombelli1987,surya2019}.
The resulting Hasse degree is bounded: $D \leq 15$.

\textbf{Phase 2: Defect detection \& generation classification}
(\texttt{skyrmion.rs}).
\emph{Causal Prisms}---complete bipartite subgraphs $K_{2,n}$ with
$n \geq 3$ mutually spacelike intermediates connecting a past pole
to a future pole---are detected via a forward-forward 2-hop search:
$O(N \cdot D^2) = O(N)$.
$K_5$ threat absorption enforces topological confinement.
Each intermediate $w_i$ carries a discrete phase
$\varphi(w_i) = \operatorname{sgn}(p_{\text{bulk}}[w_i]) \in \{-1, 0, +1\}$.
The generation of a prism is
$g(P) = |\{\text{distinct } \varphi \text{ values}\}| \in \{1, 2, 3\}$.

\textbf{Phase 3: Spectral dimension}
(\texttt{spectral.rs}).
$d_S(\sigma)$ is measured via Monte Carlo random walkers
($W = \lceil (\sigma_{\max}/\varepsilon)^2 \rceil$) on the Hasse graph.
Return probabilities are accumulated with strict integer arithmetic;
a single \texttt{f64} division is performed at the end~\cite{eichhorn2014}.

\textbf{Phase 4: Output}
(\texttt{output.rs}).
Ensemble-averaged observables with error bars, exported as CSV.

\smallskip
\textbf{Key claim:} $O(N)$ is not an optimisation---it is the
\emph{Geometry of Logic}.
Bounded Hasse degree means every topological operation is local.
This is a computational proof of physical locality as an
information-theoretic principle.

% ═══════════════════════════════════════════════════════════════════════════════
%   Section 2
% ═══════════════════════════════════════════════════════════════════════════════
\section{Results at $N = 10^6$}

Table~\ref{tab:results} summarises a single realisation (seed\,42,
runtime 63\,s).

\begin{table}[h]
\centering
\caption{Observables at $N = 10^6$ (single realisation, seed\,42).}
\label{tab:results}
\setlength{\tabcolsep}{5pt}
\begin{tabular}{@{}lc@{}}
\toprule
\textbf{Observable} & \textbf{Value} \\
\midrule
$d_S$ vacuum (global)     & 2.14 \\
$d_S$ vacuum (core)       & 3.64 \\
$d_S$ defect (core)       & 5.01 \\
\midrule
Prisms: $g{=}1$ / $g{=}2$ / $g{=}3$
                          & 5{,}731 / 26{,}228 / 3{,}565 \\
Mass: Gen\,1 / Gen\,2 / Gen\,3
                          & 4.49 / 6.51 / 7.80 \\
CPT: Gen\,1 vs Anti-1     & 4.49 vs 4.57 \\
\bottomrule
\end{tabular}
\end{table}

% ─── TikZ figure: Causal Prism ──────────────────────────────────────────────
\begin{figure}[h]
\centering
\begin{tikzpicture}[
    >=Stealth, scale=0.72, transform shape,
    pole/.style={circle, draw=black, fill=black!80, minimum size=6pt, inner sep=0pt},
    neg/.style={circle, draw=blue!70!black, fill=blue!20, minimum size=8pt, inner sep=0pt,
                font=\scriptsize\bfseries},
    zero/.style={circle, draw=black!60, fill=gray!15, minimum size=8pt, inner sep=0pt,
                 font=\scriptsize\bfseries},
    pos/.style={circle, draw=red!70!black, fill=red!20, minimum size=8pt, inner sep=0pt,
                font=\scriptsize\bfseries},
    causal/.style={->, thick, black!60}
  ]
  \node[pole, label={[font=\scriptsize]above:$u$}] (u) at (0, 2) {};
  \node[neg]  (w1) at (-2.2, 0) {$-$};
  \node[zero] (w2) at (-0.75, 0) {$0$};
  \node[pos]  (w3) at ( 0.75, 0) {$+$};
  \node[neg]  (w4) at ( 2.2, 0) {$-$};
  \node[below=1pt, font=\scriptsize] at (w1.south) {$w_1$};
  \node[below=1pt, font=\scriptsize] at (w2.south) {$w_2$};
  \node[below=1pt, font=\scriptsize] at (w3.south) {$w_3$};
  \node[below=1pt, font=\scriptsize] at (w4.south) {$w_4$};
  \node[pole, label={[font=\scriptsize]below:$v$}] (v) at (0, -2) {};
  \foreach \w in {w1,w2,w3,w4} { \draw[causal] (u) -- (\w); }
  \foreach \w in {w1,w2,w3,w4} { \draw[causal] (\w) -- (v); }
\end{tikzpicture}
\caption{A Causal Prism $K_{2,4}$.  Past pole~$u$ and future pole~$v$
  are connected through four mutually spacelike intermediates $w_i$.
  Colours encode the discrete phase
  $\varphi(w_i) \in \{{\color{blue!70!black}-},\,{\color{gray}0},\,{\color{red!70!black}+}\}$.
  This configuration has generation $g = 3$ and net phase
  $\Phi = -1$ (antimatter).}
\label{fig:prism}
\end{figure}

\textbf{Interpretation.}
The global vacuum spectral dimension $d_S = 2.14$ reflects UV dimensional
reduction, consistent with CDT~\cite{ambjorn2005} and causal set
results~\cite{carlip2017}.
At the diamond core, where boundary effects vanish, $d_S = 3.64$
---the 4D continuum is emerging.
The defect core gives $d_S = 5.01$: Causal Prisms \emph{trap} random
walkers, enhancing return probability.
This topological probability retention \emph{is} mass.

The mass hierarchy Gen\,1 $<$ Gen\,2 $<$ Gen\,3 arises because
higher-generation prisms have more complex phase structure: more distinct
$\varphi$ values produce more intermediates and stronger trapping.
CPT symmetry holds at the \textasciitilde2\% level for $M{=}1$
realisation (Gen\,1 mass 4.49 vs Anti-1 mass 4.57), expected to
sharpen with ensemble averaging.
Generation populations are naturally unequal ($g{=}2$ dominates at 26k
prisms), as dictated by the combinatorics of phase assignment.

% ═══════════════════════════════════════════════════════════════════════════════
%   Section 3
% ═══════════════════════════════════════════════════════════════════════════════
\section{Why It Works --- The Three Geometries}

The theoretical framework rests on three interlocking structures.

\textbf{Geometry of Order.}
Poisson sprinkling + transitive reduction $\to$ Hasse diagram.
The Hasse diagram is triangle-free by construction: if $u \prec w \prec v$,
the edge $u \to v$ is removed.  Therefore $K_4$ (the complete 4-clique)
cannot exist as a subgraph.
The densest bipartite structure is $K_{2,n}$---the \emph{Causal Prism}:
two causally related poles connected through $n \geq 3$ mutually
spacelike intermediates~\cite{kuratowski1930,diestel2017}.
Topological mass $= n$.

\textbf{Geometry of Matter.}
For each intermediate, $\varphi(w_i) = \operatorname{sgn}(p_{\text{bulk}}[w_i])
\in \{-1, 0, +1\}$.
The generation is $g(P) = |\{\text{distinct } \varphi\}| \in \{1, 2, 3\}$.
\textbf{Exactly three, because the sign function has three values---a theorem,
not a fit.}
The net phase $\Phi(P) = \sum_i \varphi(w_i)$ distinguishes matter
($\Phi \geq 0$) from antimatter ($\Phi < 0$).
$K_5$ threat absorption acts as topological confinement:
isolated prism components cannot exist, mirroring colour
confinement~\cite{kuratowski1930}.
The permutation symmetry $S_n$ of the intermediates
suggests emergent gauge structure.

\textbf{Geometry of Logic.}
$D \leq 15 \Rightarrow O(N) \Rightarrow$ locality.
This is not an engineering choice: it is a consequence of $\rho \approx 1$
sprinkling into the Hasse diagram of a causal diamond.
The simulation's efficiency is itself the proof~\cite{stone1936,bombelli2009}.

% ═══════════════════════════════════════════════════════════════════════════════
%   Section 4
% ═══════════════════════════════════════════════════════════════════════════════
\section{Outlook}

\begin{itemize}\setlength\itemsep{1pt}
  \item $N = 10^7$ ensemble currently running
    (target: $\varepsilon = 0.001$, 22 realisations).
  \item The $S_n$ belly symmetry of Causal Prisms is conjectured to
    generate gauge structure; proving $S_3 \to SU(3)$ from first
    principles is open.
  \item Mass ratios are observed---theoretical derivation is ongoing.
  \item All code and four pedagogic booklets (Vol.\,I, Vol.\,II,
    Kuratowski Calculus, \emph{Emma's Thought Experiments})
    are open-source at the repository above.
\end{itemize}

The Hasse diagram---already central to causal set dynamics---naturally
hosts bipartite defects whose classification reproduces the
three-generation structure of the Standard Model.
The code is available for independent verification.

% ═══════════════════════════════════════════════════════════════════════════════
%   References
% ═══════════════════════════════════════════════════════════════════════════════
\bibliographystyle{unsrt}
\bibliography{refs}

\end{document}
